% !TEX root = ../main.tex
% Resultados literales de la ejecucion sobre el corpus del curso 2026-27.
% Fuente: docs/experimentos/it38-recuperacion.md, escrito por el propio script.
% NO editar estas cifras a mano: si cambian, se vuelve a ejecutar
%   py scripts/experimento_recuperacion.py
% y se copian de ahi. Corpus de 1.499 fragmentos y 56 preguntas de dominio.
\begin{table}[htbp]
  \centering
  \small
  \caption[Recuperación sobre el conjunto de evaluación]{Recuperación del
    sistema sobre la colección completa y las 56 preguntas de dominio del
    conjunto de evaluación. Fuente: elaboración propia.}
  \label{tab:recuperacion}
  \begin{tabular}{@{}crrr@{}}
    \toprule
    \textbf{K} & \textbf{R@K} & \textbf{Techo de R@K} & \textbf{RU@K}   \\
    \midrule
    3          & 0,652        & 0,754                 & 0,906           \\
    5          & 0,789        & 0,906                 & 0,973           \\
    10         & 0,881        & 0,966                 & \textbf{0,991}  \\
    \bottomrule
  \end{tabular}
  \begin{minipage}{0.95\textwidth}
    \vspace{4pt}
    \tiny
    \textbf{Notas:} \textbf{R@K} es la exhaustividad por \emph{fragmento}
    (cuántos de los trozos de la unidad correcta aparecen entre los $K$
    primeros resultados). \textbf{RU@K} es la exhaustividad por \emph{unidad}
    (si se ha encontrado la asignatura correcta, sin penalizar que falte alguno
    de sus trozos). El \textbf{techo} es el valor más alto que R@K puede
    alcanzar sobre esta colección: hay preguntas cuyas unidades relevantes son
    más numerosas que $K$, de modo que su R@K no puede valer 1 por perfecto que
    sea el orden. Cada R@K se lee contra su techo y no contra 1; RU@K tiene
    techo 1 y se interpreta sola. El \textbf{MRR} (\textit{Mean Reciprocal
    Rank}), media del inverso de la posición del primer resultado correcto,
    es \textbf{0,926} y no depende de $K$. Las diez preguntas ajenas al dominio
    quedan fuera de estas medidas: su criterio de acierto es el contrario.
  \end{minipage}
\end{table}
